\subsection{Procesamiento y Análisis de Datos}

Una vez recolectada la información mediante los tres instrumentos
descritos en la sección anterior, el procesamiento se realizará
de forma diferenciada según la naturaleza de los datos y el objeto
de estudio del que provienen. La organización sistemática de la
información en tablas y esquemas garantizará que los resultados
sean comparables entre sí y alineados con los objetivos
específicos del proyecto.

Los datos del Instrumento N.\textsuperscript{o} 1 (entrevista
estructurada) se procesarán mediante análisis temático por
bloques. Las respuestas de los tres participantes serán transcritas
y agrupadas por área temática, identificando coincidencias y
divergencias respecto a los procesos actuales, los problemas
operativos y los requerimientos funcionales y normativos. Los
hallazgos se organizarán en una tabla de análisis que vincule
cada hallazgo con el actor que lo mencionó y con el objetivo
específico al que contribuye, facilitando la elaboración del
diagnóstico de la situación actual de Bellidental.

Los datos del Instrumento N.\textsuperscript{o} 2 (ficha de
observación de procesos) se organizarán por secciones en tablas
estructuradas. La Sección A dará lugar a un inventario de sistemas
aislados que detallará cada herramienta o archivo en uso, su tipo
y su grado de integración con los demás. La Sección B producirá un
mapa de flujo por etapa de atención, indicando el responsable, la
herramienta utilizada y el tiempo estimado por actividad. Los
hallazgos de la Sección C se clasificarán como problemas
confirmados, diferenciando los que coincidan con lo declarado en
las entrevistas de los detectados exclusivamente mediante
observación directa; esta triangulación fortalecerá el diagnóstico
final. Los resultados consolidados de ambos instrumentos
constituirán el insumo directo para el diseño del modelo de
información abordado en la fase siguiente.

Los datos del Instrumento N.\textsuperscript{o} 3 (ficha de
registro de tiempo) se analizarán con estadística descriptiva e
inferencial sobre las dos muestras de ciclos registradas en cada
fase. Los tiempos por subproceso y por ciclo completo se
registrarán en una hoja de cálculo para obtener el tiempo promedio,
la desviación estándar y los valores mínimo y máximo de cada etapa.
La comparación entre los tiempos previos y posteriores a la
implementación del sistema se realizará mediante la prueba $t$ de
Student para muestras relacionadas (diseño antes-después), con un
nivel de significancia de $\alpha = 0{,}05$. Esta prueba es
coherente con el diseño preexperimental declarado, en el que ambas
mediciones corresponden al mismo entorno operativo y al mismo
personal, antes y después de la intervención. El resultado
permitirá determinar si la reducción del tiempo promedio por ciclo
de atención es estadísticamente significativa y, por tanto, si
puede atribuirse al sistema implementado y no a la variación
aleatoria del proceso.